% ==============================================================================
% SISTEMA PÓRTICO 2.0 — DOCUMENTO MAESTRO DE MANTENIMIENTO TÉCNICO
% ==============================================================================
% El diseño, estilos, colores y paquetes se encuentran centralizados en sbc-template.sty
% Las referencias bibliográficas residen en template.bib y se formatean con sbc.bst
% ==============================================================================
\documentclass[11pt,a4paper,openany]{report}

% Carga del paquete de diseño maestro
\usepackage{sbc-template}

\begin{document}

% ------------------------------------------------------------------------------
% PORTADA EJECUTIVA MODERNA CON SEPARADORES ESTÉTICOS EN AZUL ELÉCTRICO (#1F51FF)
% ------------------------------------------------------------------------------
\begin{titlepage}
    \pagecolor{white}
    \begin{tikzpicture}[remember picture,overlay]
        % Detalle de borde minimalista
        \fill[electricblue] (current page.north west) rectangle ([xshift=5mm]current page.south west);
        \fill[midnightnavy] ([xshift=5mm]current page.north west) rectangle ([xshift=8mm]current page.south west);
        \draw[bordergray, line width=0.6pt] ([xshift=2.2cm, yshift=-2.2cm]current page.north west) rectangle ([xshift=-1.8cm, yshift=2.2cm]current page.south east);
    \end{tikzpicture}

    \vspace*{1.5cm}
    \hspace*{1.8cm}
    \begin{minipage}{0.84\textwidth}
        {\large\bfseries\color{slategray}\textls[120]{SISTEMA PÓRTICO 2.0}}\\[0.4cm]
        {\huge\bfseries\color{midnightnavy} MANUAL TÉCNICO Y GUÍA\\DE MANTENIMIENTO INTEGRAL}\\[0.4cm]
        {\large\color{slategray} Especificación de Arquitectura Híbrida Dual, Flujos de Métodos,\\Persistencia Relacional y Protocolos de Extensión de Software}\\[0.6cm]

        % Separador estético en azul eléctrico (#1F51FF)
        \electricrule
        \vspace*{0.6cm}

        \begin{minipage}[t]{0.48\textwidth}
            \textbf{\color{midnightnavy}Gobernanza del Sistema:}\\[0.2cm]
            \textbf{Plataforma:} Sistema Pórtico\\
            \textbf{Arquitectura:} Híbrida Dual (Web / Desktop)\\
            \textbf{Backend Web:} PHP 8.2+ / Apache / PDO\\
            \textbf{Host Nativo:} .NET 10 / C\# / Edge WebView2\\
            \textbf{Persistencia:} MySQL 8.0 (Motor InnoDB)
        \end{minipage}
        \hfill
        \begin{minipage}[t]{0.48\textwidth}
            \textbf{\color{midnightnavy}Metadatos del Documento:}\\[0.2cm]
            \textbf{Autoría:} Tiago Pereira C.\\
            \textbf{Rol:} Ingeniería de Software y DevOps\\
            \textbf{Clasificación:} Documentación Técnica Confidencial\\
            \textbf{Fecha de Edición:} Septiembre de 2026\\
            \textbf{Estado de Entrega:} Versión 2.0 Maestro
        \end{minipage}

        \vfill
        \vspace*{3.5cm}

        \begin{modernbox}{PROPÓSITO Y OBLIGATORIEDAD DE ESTA GUÍA}
            Este manual técnico constituye la fuente canónica de arquitectura, directrices de diseño y gobernanza de código para futuros ciclos de mantenimiento evolutivo y correctivo del sistema Pórtico. Cualquier modificación que rompa la paridad estricta entre la API REST (PHP) y el Host de Escritorio (.NET 10) compromete la integridad operativa de la plataforma.
        \end{modernbox}
    \end{minipage}
\end{titlepage}
\nopagecolor

% ------------------------------------------------------------------------------
% TABLA DE CONTENIDOS, FIGURAS Y TABLAS
% ------------------------------------------------------------------------------
\tableofcontents
\newpage
\listoffigures
\listoftables
\newpage

% ==============================================================================
% CAPÍTULO 1: INTRODUCCIÓN Y ARQUITECTURA GENERAL
% ==============================================================================
\chapter{Fundamentos y Arquitectura del Sistema}

\section{Propósito del Sistema y Objetivos de Mantenimiento}

El sistema \textbf{Pórtico} es una plataforma empresarial diseñada para centralizar y optimizar la administración integral de proyectos constructivos, cuadrillas de personal operativo, parque de maquinaria pesada, suministros energéticos y trazabilidad de contratos legales.

A lo largo del ciclo de vida de una obra civil, surgen múltiples vectores de inconsistencia: desfasaje en los partes diarios de asistencia, demoras en la detección de certificados técnicos vencidos, liquidaciones imprecisas de tareas contractuales y falta de visibilidad en el consumo de combustible por máquina. Pórtico resuelve estos desafíos a través de un modelo de datos unificado y una interfaz gráfica responsiva.

\electricsep

\section{Arquitectura Global: El Modelo Híbrido Dual}

La característica más relevante del sistema Pórtico reside en su \textbf{Arquitectura Híbrida Dual Simétrica}: la misma capa de presentación visual (\texttt{web-ui/}), construida con estándares web abiertos (HTML5, CSS modular y JavaScript Vanilla), es ejecutada en dos runtimes con entornos de ejecución radicalmente distintos pero con idéntica lógica de negocio y persistencia en una base de datos MySQL centralizada.

\begin{figure}[H]
    \centering
    \noindent\resizebox{\linewidth}{!}{%
    \begin{tikzpicture}[
        node distance=1.4cm and 1.5cm,
        box/.style={rectangle, rounded corners=1.5mm, draw=bordergray, line width=0.8pt, fill=white, inner sep=7pt, align=center, font=\small},
        ui/.style={box, draw=electricblue, fill=cardbg, font=\bfseries\small, text width=13.5cm, align=center},
        web/.style={box, draw=bordergray, fill=white},
        win/.style={box, draw=bordergray, fill=white},
        db/.style={cylinder, cylinder uses custom fill, cylinder end fill=cardbg, cylinder body fill=white, draw=midnightnavy, line width=1.2pt, aspect=0.25, shape border rotate=90, inner sep=6pt, align=center, font=\bfseries\small},
        arrow/.style={-{Stealth[scale=1.0]}, line width=1.1pt, color=electricblue}
    ]
        % Nodo Frontend
        \node[ui, minimum width=13.5cm] (frontend) {
            \textbf{\color{midnightnavy}CAPA DE PRESENTACIÓN COMPARTIDA (\texttt{web-ui/})}\\[0.1cm]
            \footnotesize Vistas HTML5 semánticas \textbar\ Estilos CSS con Tokens \textbar\ Controladores JS con autodetección de entorno
        };

        % Runtimes
        \node[web, below left=1.8cm and -0.6cm of frontend, text width=6.2cm] (runtimeWeb) {
            \textbf{\color{midnightnavy}RUNTIME 1: SERVIDOR WEB (PHP 8+)}\\[0.1cm]
            \footnotesize $\bullet$ Apache / Nginx / PHP CLI Server\\
            $\bullet$ Endpoints RESTful (\texttt{api/*.php})\\
            $\bullet$ Filtros Anti-CSRF \& Control de Sesión\\
            $\bullet$ Driver PDO MySQL (Prepared Statements)
        };

        \node[win, below right=1.8cm and -0.6cm of frontend, text width=6.2cm] (runtimeWin) {
            \textbf{\color{midnightnavy}RUNTIME 2: ESCRITORIO (.NET 10)}\\[0.1cm]
            \footnotesize $\bullet$ Host Windows Forms (\texttt{PorticoDesktop.exe})\\
            $\bullet$ Motor Microsoft Edge WebView2\\
            $\bullet$ Esquema Virtual (\texttt{https://portico.desktop/*})\\
            $\bullet$ Canal IPC Bidireccional (\texttt{postMessage})\\
            $\bullet$ Driver Directo ADO.NET (\texttt{MySqlConnector})
        };

        % Base de Datos
        \node[db, below=2.5cm of frontend, text width=9.2cm] (database) {
            \color{midnightnavy}MOTOR DE PERSISTENCIA RELACIONAL: MySQL 8.0 (InnoDB)\\[0.1cm]
            \footnotesize 18 Tablas Relacionales \textbar\ Esquema DDL Maestro: \texttt{schema.sql} \textbar\ Transacciones ACID
        };

        % Conectores
        \draw[arrow] (frontend.south -| runtimeWeb.north) -- node[left, font=\scriptsize\color{slategray}] {HTTP(S) / \texttt{fetch()}} (runtimeWeb.north);
        \draw[arrow] (frontend.south -| runtimeWin.north) -- node[right, font=\scriptsize\color{slategray}] {IPC \texttt{window.chrome.webview}} (runtimeWin.north);

        \draw[arrow] (runtimeWeb.south) -- node[below left, font=\scriptsize\color{slategray}] {TCP 3306 (PDO)} (database.north -| runtimeWeb.south);
        \draw[arrow] (runtimeWin.south) -- node[below right, font=\scriptsize\color{slategray}] {TCP 3306 (ADO.NET)} (database.north -| runtimeWin.south);
    \end{tikzpicture}%
    }
    \caption{Diagrama de Arquitectura Híbrida Dual: Separación de Runtimes y Persistencia.}
    \label{fig:arquitectura_tikz}
\end{figure}

\begin{figure}[H]
    \centering
    \begin{modernbox}{CAPTURA DEL DIAGRAMA COMPLEMENTARIO (MERMAID)}
        \centering\vspace{0.8cm}
        \textbf{\color{slategray}[CAPTURA MERMAID OPCIONAL: fig\_arquitectura\_hibrida.png]}\\[0.15cm]
        \footnotesize Si desea insertar la captura generada en Markdown, ubique la imagen en la carpeta \texttt{img/} y descomente la línea de inclusión.
        \vspace{0.8cm}
    \end{modernbox}
    \caption{Representación complementaria de la Arquitectura Híbrida del Sistema.}
    \label{fig:arquitectura_mermaid}
\end{figure}

\electricsep

\section{Justificación Técnica y Beneficios del Modelo}

La coexistencia de dos runtimes proporciona ventajas operativas insustituibles en la industria de la construcción:
\begin{enumerate}
    \item \textbf{Cero Latencia en Oficinas Centrales (Modo Escritorio):} Los supervisores y directores de proyecto disponen de una aplicación nativa en Windows (.NET 10) que interactúa con la base de datos a través de sockets TCP directos y transacciones compiladas, brindando tiempos de respuesta instantáneos sin overhead de capa web intermedia.
    \item \textbf{Movilidad y Operatividad en Campo (Modo Web):} Los capataces en pie de obra pueden acceder desde dispositivos móviles, tablets o portátiles a través de cualquier navegador conectándose al servidor PHP (\texttt{api/}), registrando asistencias y combustible en tiempo real.
    \item \textbf{Mantenibilidad Centralizada del Frontend:} Todo cambio en la interfaz gráfica (diseño visual, validaciones de formularios, interacción DOM) se efectúa una única vez en la carpeta \texttt{web-ui/}. Dado que el proyecto .NET incrusta dichos archivos como recursos binarios en tiempo de compilación (\texttt{EmbeddedResource}), no existe duplicación de código de frontend.
\end{enumerate}

% ==============================================================================
% CAPÍTULO 2: MODELO HÍBRIDO DUAL
% ==============================================================================
\chapter{El Modelo Híbrido Dual: Web (PHP 8) vs. Escritorio (.NET 10)}

\section{Runtime Web: Pila Tecnológica y Flujo de Ejecución}

En el entorno Web, el sistema opera bajo una arquitectura de microcontroladores orientada a recursos:
\begin{itemize}
    \item \textbf{Servidor Web:} Compatible con Apache, Nginx o el servidor embebido de desarrollo PHP (\texttt{php -S localhost:8000}).
    \item \textbf{Control de Sesión Reforzado:} Implementado en \texttt{api/config/session.php}, establece parámetros de cookies de extrema seguridad (\texttt{SameSite=Strict}, \texttt{HttpOnly=true} y \texttt{Secure} condicional si opera bajo HTTPS).
    \item \textbf{Protección contra Ataques CSRF:} Cada petición de mutación (\texttt{POST}, \texttt{PUT}, \texttt{DELETE}) valida que la cabecera enviada por el cliente (\texttt{X-CSRF-Token}) coincida con el token criptográfico persistido en la sesión del servidor (\texttt{\$\_SESSION['csrf\_token']}) empleando la función segura \texttt{hash\_equals()}.
    \item \textbf{Detección de Desbordamiento de Carga (\texttt{post\_max\_size}):} La función \texttt{verificarLimitePost()} en \texttt{api/config/utils.php} detecta cuando un usuario sube un contrato o plano que excede el límite fijado en el servidor, impidiendo que PHP descarte silenciosamente el payload y arrojando una excepción HTTP 400 legible.
\end{itemize}

\begin{lstlisting}[language=PHP, caption={Patrón de Conexión Singleton Seguro en PHP (api/config/db.php)}]
function conectar(): PDO {
    static $pdo = null;
    if ($pdo !== null) return $pdo;

    // Carga de variables de entorno desde .env
    $host   = $_ENV['DB_HOST']     ?? 'localhost';
    $port   = $_ENV['DB_PORT']     ?? '3306';
    $dbname = $_ENV['DB_NAME']     ?? 'portico';
    $user   = $_ENV['DB_USER']     ?? '';
    $pass   = $_ENV['DB_PASSWORD'] ?? '';

    $dsn = "mysql:host={$host};port={$port};dbname={$dbname};charset=utf8mb4";
    $pdo = new PDO($dsn, $user, $pass, [
        PDO::ATTR_ERRMODE            => PDO::ERRMODE_EXCEPTION,
        PDO::ATTR_DEFAULT_FETCH_MODE => PDO::FETCH_ASSOC,
        PDO::ATTR_EMULATE_PREPARES   => false, // Inmunidad a SQLi
        PDO::ATTR_TIMEOUT            => 5,
    ]);
    return $pdo;
}
\end{lstlisting}

\electricsep

\section{Runtime Desktop: .NET 10, WebView2 e IPC Nativo}

La aplicación de escritorio prescinde totalmente de servidores web locales (como Kestrel o Apache local). Su arquitectura se fundamenta en:
\begin{enumerate}
    \item \textbf{Host WinForms y Control WebView2:} La ventana principal (\texttt{MainWindow.cs}) aloja el control \texttt{Microsoft.Web.WebView2.WinForms.WebView2}, inicializando el runtime de Chromium.
    \item \textbf{Filtro de Recursos Embebidos:} Se establece una regla de resolución virtual:
    \begin{lstlisting}[language={[Sharp]C}]
_webView.CoreWebView2.AddWebResourceRequestedFilter(
    "https://portico.desktop/*", 
    CoreWebView2WebResourceContext.All
);
    \end{lstlisting}
    Cuando el navegador interno solicita \texttt{https://portico.desktop/Obras.html}, el evento \texttt{OnWebResourceRequested} busca el recurso incrustado en el ensamblado bajo la convención \texttt{webui.Obras.html} y devuelve un flujo de lectura en memoria (\texttt{Stream}) con el tipo MIME correspondiente.
    \item \textbf{Bucle de Mensajes IPC (Inter-Process Communication):} Para ejecutar operaciones contra la base de datos, el código JavaScript no realiza llamadas \texttt{fetch()}, sino que emite mensajes estructurados:
    \begin{lstlisting}[language=javascript]
window.chrome.webview.postMessage(JSON.stringify({
    type: 'obras_guardar',
    requestId: 'req_1726000000',
    nombre: 'Ampliacion Edificio Norte',
    numero_contrata: 'OB-2026-042'
}));
    \end{lstlisting}
    El método \texttt{OnWebMessageReceived} en C\# intercepta el JSON, valida la sesión en memoria, ejecuta la transacción contra MySQL mediante \texttt{MySqlConnector} y responde a través de \texttt{PostWebMessageAsString()} emparejando el mismo \texttt{requestId}.
\end{enumerate}

\begin{alertbox}{PARIDAD ESTRICTA ENTRE API Y C\#}
Cualquier modificación o añadido en una ruta de la API PHP (\texttt{api/*.php}) \textbf{DEBE CONTAR OBLIGATORIAMENTE} con su homólogo despachador en el bloque \texttt{switch (messageType)} de \texttt{MainWindow.cs}. Omitir esta regla inutilizará la vista correspondiente en el ejecutable de escritorio.
\end{alertbox}

% ==============================================================================
% CAPÍTULO 3: MAPA MENTAL Y TAXONOMÍA
% ==============================================================================
\chapter{Mapa Mental y Taxonomía del Ecosistema Pórtico}

\section{Taxonomía Jerárquica del Repositorio}

El mapa mental de la Figura \ref{fig:mapamental_tikz} ilustra la descomposición funcional del sistema, estructurado en cuatro ramas maestras: Capa de Presentación, Runtime Backend Web, Runtime Desktop y Motor de Persistencia Relacional.

\begin{figure}[H]
    \centering
    \noindent\resizebox{0.88\linewidth}{!}{%
    \begin{tikzpicture}[
        mindmap,
        grow cyclic,
        every node/.style={concept},
        root concept/.append style={concept color=midnightnavy, text=white, font=\bfseries\small, minimum size=2.6cm},
        level 1 concept/.append style={font=\bfseries\scriptsize, minimum size=2.1cm, text width=1.9cm},
        level 2 concept/.append style={font=\tiny, minimum size=1.4cm, text width=1.5cm},
        p1/.style={concept color=electricblue, text=white},
        p2/.style={concept color=slategray, text=white},
        p3/.style={concept color=charcoal, text=white},
        p4/.style={concept color=midnightnavy!80!black, text=white}
    ]
        \node[root concept] {SISTEMA PÓRTICO}
            child[concept color=electricblue] {
                node[p1] {Capa Visual web-ui}
                child { node {HTML5 Semántico} }
                child { node {CSS Tokens} }
                child { node {JS Vanilla} }
            }
            child[concept color=slategray] {
                node[p2] {Backend PHP 8}
                child { node {Endpoints REST} }
                child { node {Anti-CSRF} }
                child { node {PDO MySQL} }
            }
            child[concept color=charcoal] {
                node[p3] {Host Desktop C\#}
                child { node {WinForms 10} }
                child { node {WebView2 IPC} }
                child { node {MySql\\Connector} }
            }
            child[concept color=midnightnavy!80!black] {
                node[p4] {Persistencia MySQL}
                child { node {18 Tablas} }
                child { node {InnoDB ACID} }
                child { node {schema.sql} }
            };
    \end{tikzpicture}%
    }
    \caption{Mapa Mental: Descomposición jerárquica de componentes del sistema.}
    \label{fig:mapamental_tikz}
\end{figure}

\begin{figure}[H]
    \centering
    \begin{modernbox}{CAPTURA DEL MAPA MENTAL COMPLEMENTARIO (MERMAID)}
        \centering\vspace{0.8cm}
        \textbf{\color{slategray}[CAPTURA MERMAID OPCIONAL: fig\_mapa\_mental.png]}\\[0.15cm]
        \footnotesize Inserte aquí la captura del mapa mental generado en el documento Markdown.
        \vspace{0.8cm}
    \end{modernbox}
    \caption{Representación complementaria del Mapa Mental del Sistema.}
    \label{fig:mapa_mermaid}
\end{figure}

% ==============================================================================
% CAPÍTULO 4: MODELO DE BASE DE DATOS
% ==============================================================================
\chapter{Catálogo y Diccionario Relacional de Base de Datos}

\section{Diagrama Entidad-Relación Consolidado}

El modelo relacional está compuesto por 18 tablas en tercera forma normal (3NF), optimizadas para operaciones concurrentes mediante el motor transaccional \textbf{InnoDB} con soporte nativo para caracteres UTF-8 extendidos (\texttt{utf8mb4\_unicode\_ci}).

\begin{figure}[H]
    \centering
    \begin{modernbox}{DIAGRAMA ENTIDAD-RELACIÓN COMPLEMENTARIO}
        \centering\vspace{1.2cm}
        \textbf{\color{slategray}[DIAGRAMA ERD COMPLETO: fig\_erd\_completo.png]}\\[0.15cm]
        \footnotesize Inserte aquí la captura del Diagrama Entidad-Relación exportado desde Markdown o herramientas de diseño.
        \vspace{1.2cm}
    \end{modernbox}
    \caption{Diagrama Entidad-Relación Consolidado (18 Tablas Relacionales).}
    \label{fig:erd_placeholder}
\end{figure}

\electricsep

\section{Diccionario Técnico de Tablas y Restricciones}

A continuación se detalla la especificación de cada tabla del esquema maestro \texttt{schema.sql}, configurada con anchos balanceados y ajuste de línea automático para evitar desbordamientos:

% 1. usuarios
\subsection{Tabla \texttt{usuarios}}
Almacena las credenciales y roles del personal con acceso al sistema.
\begin{table}[H]
\centering\small
\renewcommand{\arraystretch}{1.3}
\begin{tabularx}{\textwidth}{>{\ttfamily\bfseries\color{midnightnavy}\raggedright\arraybackslash}p{2.7cm} >{\ttfamily\raggedright\arraybackslash}p{3.5cm} >{\raggedright\arraybackslash}p{2.2cm} Y}
\rowcolor{slate!8}
\textbf{Campo} & \textbf{Tipo de Dato} & \textbf{Nulo / Default} & \textbf{Descripción y Reglas de Integridad} \\
\specialrule{1pt}{0pt}{2pt}
id\_usuario & INT AUTO\_INCREMENT & NOT NULL, PK & Identificador único del usuario. \\
nombre & VARCHAR(100) & NOT NULL & Nombre y apellido completo. \\
usuario & VARCHAR(50) & NOT NULL, UNIQUE & Nombre de usuario único para inicio de sesión. \\
password\_hash & VARCHAR(255) & NOT NULL & Hash criptográfico BCrypt (\$2y\$ / \$2a\$). \\
correo & VARCHAR(150) & NULL, UNIQUE & Correo electrónico institucional. \\
rol & ENUM('Admin', 'Capataz') & NOT NULL & Control de acceso por roles RBAC. \\
activo & BOOLEAN & DEFAULT 1 & Borrado lógico de cuentas (1: Activo, 0: Baja). \\
\bottomrule
\end{tabularx}
\caption{Estructura técnica de la tabla \texttt{usuarios}.}
\label{tab:usuarios}
\end{table}

% 2. intentos_login
\subsection{Tabla \texttt{intentos\_login}}
Auditoría y mitigación de ataques de fuerza bruta en el inicio de sesión.
\begin{table}[H]
\centering\small
\renewcommand{\arraystretch}{1.3}
\begin{tabularx}{\textwidth}{>{\ttfamily\bfseries\color{midnightnavy}\raggedright\arraybackslash}p{2.7cm} >{\ttfamily\raggedright\arraybackslash}p{3.5cm} >{\raggedright\arraybackslash}p{2.2cm} Y}
\rowcolor{slate!8}
\textbf{Campo} & \textbf{Tipo de Dato} & \textbf{Nulo / Default} & \textbf{Descripción y Reglas de Integridad} \\
\specialrule{1pt}{0pt}{2pt}
id & INT AUTO\_INCREMENT & NOT NULL, PK & Identificador del intento de autenticación. \\
username & VARCHAR(50) & NOT NULL & Nombre de usuario ingresado en el intento. \\
ip\_address & VARCHAR(45) & NOT NULL & Dirección IPv4 o IPv6 del cliente. \\
fecha & DATETIME & NOT NULL & Marca temporal del evento de acceso. \\
exitoso & BOOLEAN & DEFAULT 0 & Resultado de la validación (1: Válido, 0: Fallido). \\
\bottomrule
\end{tabularx}
\caption{Estructura técnica de la tabla \texttt{intentos\_login}.}
\label{tab:intentos_login}
\end{table}

% 3. auditoria_logs
\subsection{Tabla \texttt{auditoria\_logs}}
Pista forense inmutable de todas las mutaciones críticas en el sistema.
\begin{table}[H]
\centering\small
\renewcommand{\arraystretch}{1.3}
\begin{tabularx}{\textwidth}{>{\ttfamily\bfseries\color{midnightnavy}\raggedright\arraybackslash}p{2.7cm} >{\ttfamily\raggedright\arraybackslash}p{3.5cm} >{\raggedright\arraybackslash}p{2.2cm} Y}
\rowcolor{slate!8}
\textbf{Campo} & \textbf{Tipo de Dato} & \textbf{Nulo / Default} & \textbf{Descripción y Reglas de Integridad} \\
\specialrule{1pt}{0pt}{2pt}
id\_log & INT AUTO\_INCREMENT & NOT NULL, PK & Identificador único del evento de auditoría. \\
id\_usuario & INT & NULL, FK & Usuario ejecutor (ON DELETE SET NULL). \\
usuario & VARCHAR(50) & NULL & Nombre de usuario en el momento de la acción. \\
rol & VARCHAR(20) & NULL & Rol administrativo del operador. \\
accion & VARCHAR(20) & NOT NULL & Tipo de operación (LOGIN, CREAR, EDITAR). \\
entidad & VARCHAR(50) & NOT NULL & Módulo afectado (obras, obreros, contratos). \\
entidad\_id & INT & NULL & ID del registro intervenido. \\
detalle\_json & TEXT & NULL & Snapshot JSON con el estado previo y posterior. \\
ip\_address & VARCHAR(45) & NULL & Dirección IP de la terminal emisora. \\
created\_at & DATETIME & NOT NULL & Marca de tiempo del registro. \\
\bottomrule
\end{tabularx}
\caption{Estructura técnica de la tabla \texttt{auditoria\_logs}.}
\label{tab:auditoria_logs}
\end{table}

% 4. obras
\subsection{Tabla \texttt{obras}}
Entidad principal del dominio que modela los proyectos constructivos.
\begin{table}[H]
\centering\small
\renewcommand{\arraystretch}{1.3}
\begin{tabularx}{\textwidth}{>{\ttfamily\bfseries\color{midnightnavy}\raggedright\arraybackslash}p{2.7cm} >{\ttfamily\raggedright\arraybackslash}p{3.5cm} >{\raggedright\arraybackslash}p{2.2cm} Y}
\rowcolor{slate!8}
\textbf{Campo} & \textbf{Tipo de Dato} & \textbf{Nulo / Default} & \textbf{Descripción y Reglas de Integridad} \\
\specialrule{1pt}{0pt}{2pt}
id\_obra & INT AUTO\_INCREMENT & NOT NULL, PK & Identificador primario de la obra civil. \\
numero\_contrata & VARCHAR(50) & NOT NULL, UNIQUE & Código administrativo oficial de la contrata. \\
nombre & VARCHAR(150) & NOT NULL & Denominación formal del proyecto. \\
direccion & VARCHAR(200) & NULL & Emplazamiento geográfico de la obra. \\
descripcion & TEXT & NULL & Memoria descriptiva y alcance técnico. \\
fecha\_inicio & DATE & NULL & Fecha de apertura de actividades de obra. \\
fecha\_fin & DATE & NULL & Fecha prevista o efectiva de finalización. \\
nombre\_cliente & VARCHAR(150) & NOT NULL & Razón social o titular contratante. \\
telefono\_cliente & VARCHAR(30) & NULL & Teléfono de contacto institucional. \\
activo & BOOLEAN & DEFAULT 1 & Flag de operatividad (1: En ejecución, 0: Concluida). \\
\bottomrule
\end{tabularx}
\caption{Estructura técnica de la tabla \texttt{obras}.}
\label{tab:obras}
\end{table}

% 5. contratos
\subsection{Tabla \texttt{contratos}}
Almacena los documentos de contrata legal y control financiero por obra.
\begin{table}[H]
\centering\small
\renewcommand{\arraystretch}{1.3}
\begin{tabularx}{\textwidth}{>{\ttfamily\bfseries\color{midnightnavy}\raggedright\arraybackslash}p{2.7cm} >{\ttfamily\raggedright\arraybackslash}p{3.5cm} >{\raggedright\arraybackslash}p{2.2cm} Y}
\rowcolor{slate!8}
\textbf{Campo} & \textbf{Tipo de Dato} & \textbf{Nulo / Default} & \textbf{Descripción y Reglas de Integridad} \\
\specialrule{1pt}{0pt}{2pt}
id\_contrato & INT AUTO\_INCREMENT & NOT NULL, PK & Identificador del contrato. \\
id\_obra & INT & NOT NULL, FK & Referencia a \texttt{obras} (ON DELETE CASCADE). \\
archivo & LONGBLOB & NOT NULL & Binario del contrato en PDF/Scan (hasta 10 MB). \\
nombre\_archivo & VARCHAR(255) & NULL & Nombre original del fichero cargado. \\
fecha\_subida & DATE & NOT NULL & Fecha de ingreso documental al sistema. \\
id\_contrato\_origen & INT & NULL, FK & Contrato predecesor en adendas (SET NULL). \\
estado & ENUM('Activo', 'Cerrado') & DEFAULT 'Activo' & Estado del contrato en su ciclo de vida. \\
fecha\_cierre & DATE & NULL & Fecha de liquidación formal. \\
monto\_total & DECIMAL(12,2) & DEFAULT 0.00 & Presupuesto total del contrato. \\
monto\_liquidado & DECIMAL(12,2) & DEFAULT 0.00 & Suma de tareas efectivamente completadas. \\
importe\_final & DECIMAL(12,2) & DEFAULT 0.00 & Saldo o liquidación final consolidada. \\
motivo\_cierre & ENUM(...) & NULL & 'Nuevo contrato' o 'Finalizacion de obra'. \\
\bottomrule
\end{tabularx}
\caption{Estructura técnica de la tabla \texttt{contratos}.}
\label{tab:contratos}
\end{table}

% 6. contrato_tareas
\subsection{Tabla \texttt{contrato\_tareas}}
Desglose presupuestario de tareas o partidas contractuales.
\begin{table}[H]
\centering\small
\renewcommand{\arraystretch}{1.3}
\begin{tabularx}{\textwidth}{>{\ttfamily\bfseries\color{midnightnavy}\raggedright\arraybackslash}p{2.7cm} >{\ttfamily\raggedright\arraybackslash}p{3.5cm} >{\raggedright\arraybackslash}p{2.2cm} Y}
\rowcolor{slate!8}
\textbf{Campo} & \textbf{Tipo de Dato} & \textbf{Nulo / Default} & \textbf{Descripción y Reglas de Integridad} \\
\specialrule{1pt}{0pt}{2pt}
id\_tarea & INT AUTO\_INCREMENT & NOT NULL, PK & Identificador único de la partida o tarea. \\
id\_contrato & INT & NOT NULL, FK & Contrato asociado (ON DELETE CASCADE). \\
id\_tarea\_origen & INT & NULL & ID de la tarea original si proviene de una adenda. \\
descripcion & VARCHAR(255) & NOT NULL & Detalle técnico de la tarea o hito. \\
importe & DECIMAL(12,2) & DEFAULT 0.00 & Valor asignado en el presupuesto. \\
estado & ENUM('Pendiente', 'Completada') & DEFAULT 'Pendiente' & Avance físico del ítem. \\
fecha\_completada & DATE & NULL & Fecha de validación por el capataz. \\
\bottomrule
\end{tabularx}
\caption{Estructura técnica de la tabla \texttt{contrato\_tareas}.}
\label{tab:contrato_tareas}
\end{table}

% 7. obreros
\subsection{Tabla \texttt{obreros}}
Padrón de personal operativo y cuadrillas de construcción.
\begin{table}[H]
\centering\small
\renewcommand{\arraystretch}{1.3}
\begin{tabularx}{\textwidth}{>{\ttfamily\bfseries\color{midnightnavy}\raggedright\arraybackslash}p{2.7cm} >{\ttfamily\raggedright\arraybackslash}p{3.5cm} >{\raggedright\arraybackslash}p{2.2cm} Y}
\rowcolor{slate!8}
\textbf{Campo} & \textbf{Tipo de Dato} & \textbf{Nulo / Default} & \textbf{Descripción y Reglas de Integridad} \\
\specialrule{1pt}{0pt}{2pt}
id\_obrero & INT AUTO\_INCREMENT & NOT NULL, PK & Identificador del trabajador. \\
nombre & VARCHAR(150) & NOT NULL & Nombres del operario. \\
apellido & VARCHAR(100) & NULL & Apellidos del operario. \\
documento & VARCHAR(30) & NOT NULL, UNIQUE & DNI, CUIT o documento de identidad. \\
telefono & VARCHAR(30) & NULL & Teléfono personal de contacto. \\
fecha\_contratacion & DATE & NULL & Fecha de alta laboral. \\
fecha\_fin & DATE & NULL & Fecha de baja o cese laboral. \\
cargo & ENUM(...) & NOT NULL & Oficio (Albañil, Capataz, Electricista, Peón, etc.). \\
activo & BOOLEAN & DEFAULT 1 & Flag de operatividad (1: En plantilla, 0: Baja). \\
\bottomrule
\end{tabularx}
\caption{Estructura técnica de la tabla \texttt{obreros}.}
\label{tab:obreros}
\end{table}

% 8. contrato_obrero
\subsection{Tabla \texttt{contrato\_obrero}}
Legajos laborales y pólizas con control de vencimiento.
\begin{table}[H]
\centering\small
\renewcommand{\arraystretch}{1.3}
\begin{tabularx}{\textwidth}{>{\ttfamily\bfseries\color{midnightnavy}\raggedright\arraybackslash}p{2.7cm} >{\ttfamily\raggedright\arraybackslash}p{3.5cm} >{\raggedright\arraybackslash}p{2.2cm} Y}
\rowcolor{slate!8}
\textbf{Campo} & \textbf{Tipo de Dato} & \textbf{Nulo / Default} & \textbf{Descripción y Reglas de Integridad} \\
\specialrule{1pt}{0pt}{2pt}
id\_contrato\_obrero & INT AUTO\_INCREMENT & NOT NULL, PK & Identificador del legajo digital. \\
archivo & LONGBLOB & NOT NULL & Contenido binario del contrato digitalizado. \\
nombre\_archivo & VARCHAR(255) & NULL & Nombre original del archivo PDF/imagen. \\
id\_obrero & INT & NOT NULL, FK & Obrero titular (ON DELETE CASCADE). \\
fecha\_vencimiento & DATE & NULL & Fecha de caducidad para disparo de alertas. \\
\bottomrule
\end{tabularx}
\caption{Estructura técnica de la tabla \texttt{contrato\_obrero}.}
\label{tab:contrato_obrero}
\end{table}

% 9. maquinaria
\subsection{Tabla \texttt{maquinaria}}
Catálogo maestro de vehículos, maquinaria pesada y equipos de obra.
\begin{table}[H]
\centering\small
\renewcommand{\arraystretch}{1.3}
\begin{tabularx}{\textwidth}{>{\ttfamily\bfseries\color{midnightnavy}\raggedright\arraybackslash}p{2.7cm} >{\ttfamily\raggedright\arraybackslash}p{3.5cm} >{\raggedright\arraybackslash}p{2.2cm} Y}
\rowcolor{slate!8}
\textbf{Campo} & \textbf{Tipo de Dato} & \textbf{Nulo / Default} & \textbf{Descripción y Reglas de Integridad} \\
\specialrule{1pt}{0pt}{2pt}
id\_maquinaria & INT AUTO\_INCREMENT & NOT NULL, PK & Identificador del equipo. \\
nombre & VARCHAR(150) & NOT NULL & Descripción o denominación de la máquina. \\
marca & VARCHAR(100) & NULL & Fabricante o modelo técnico del equipo. \\
\bottomrule
\end{tabularx}
\caption{Estructura técnica de la tabla \texttt{maquinaria}.}
\label{tab:maquinaria}
\end{table}

% 10. certificado
\subsection{Tabla \texttt{certificado}}
Pólizas técnicas, seguros de cobertura y habilitaciones de maquinaria.
\begin{table}[H]
\centering\small
\renewcommand{\arraystretch}{1.3}
\begin{tabularx}{\textwidth}{>{\ttfamily\bfseries\color{midnightnavy}\raggedright\arraybackslash}p{2.7cm} >{\ttfamily\raggedright\arraybackslash}p{3.5cm} >{\raggedright\arraybackslash}p{2.2cm} Y}
\rowcolor{slate!8}
\textbf{Campo} & \textbf{Tipo de Dato} & \textbf{Nulo / Default} & \textbf{Descripción y Reglas de Integridad} \\
\specialrule{1pt}{0pt}{2pt}
id\_certificado & INT AUTO\_INCREMENT & NOT NULL, PK & Identificador del certificado técnico. \\
archivo & LONGBLOB & NOT NULL & Documento digitalizado de la certificación. \\
nombre\_archivo & VARCHAR(255) & NULL & Nombre original del comprobante. \\
id\_maquinaria & INT & NOT NULL, FK & Equipo vinculado (ON DELETE CASCADE). \\
fecha\_vencimiento & DATE & NULL & Fecha de caducidad para disparo de alertas. \\
\bottomrule
\end{tabularx}
\caption{Estructura técnica de la tabla \texttt{certificado}.}
\label{tab:certificado}
\end{table}

% 11. obra_maquinaria
\subsection{Tabla \texttt{obra\_maquinaria}}
Historial de asignación y destino de maquinaria pesada en proyectos.
\begin{table}[H]
\centering\small
\renewcommand{\arraystretch}{1.3}
\begin{tabularx}{\textwidth}{>{\ttfamily\bfseries\color{midnightnavy}\raggedright\arraybackslash}p{2.7cm} >{\ttfamily\raggedright\arraybackslash}p{3.5cm} >{\raggedright\arraybackslash}p{2.2cm} Y}
\rowcolor{slate!8}
\textbf{Campo} & \textbf{Tipo de Dato} & \textbf{Nulo / Default} & \textbf{Descripción y Reglas de Integridad} \\
\specialrule{1pt}{0pt}{2pt}
id\_obra\_maquinaria & INT AUTO\_INCREMENT & NOT NULL, PK & Identificador de la asignación. \\
id\_obra & INT & NOT NULL, FK & Obra receptora (ON DELETE CASCADE). \\
id\_maquinaria & INT & NOT NULL, FK & Maquinaria asignada (ON DELETE CASCADE). \\
fecha\_asignacion & DATE & NULL & Fecha de ingreso al predio constructivo. \\
fecha\_retiro & DATE & NULL & Fecha de desmovilización. \\
\bottomrule
\end{tabularx}
\caption{Estructura técnica de la tabla \texttt{obra\_maquinaria}.}
\label{tab:obra_maquinaria}
\end{table}

% 12. asistencia_maquinaria
\subsection{Tabla \texttt{asistencia\_maquinaria}}
Control diario de entradas y salidas de maquinaria por obra.
\begin{table}[H]
\centering\small
\renewcommand{\arraystretch}{1.3}
\begin{tabularx}{\textwidth}{>{\ttfamily\bfseries\color{midnightnavy}\raggedright\arraybackslash}p{2.7cm} >{\ttfamily\raggedright\arraybackslash}p{3.5cm} >{\raggedright\arraybackslash}p{2.2cm} Y}
\rowcolor{slate!8}
\textbf{Campo} & \textbf{Tipo de Dato} & \textbf{Nulo / Default} & \textbf{Descripción y Reglas de Integridad} \\
\specialrule{1pt}{0pt}{2pt}
id & INT AUTO\_INCREMENT & NOT NULL, PK & Identificador del movimiento de máquina. \\
id\_obra & INT & NOT NULL, FK & Obra asignada (ON DELETE CASCADE). \\
id\_maquinaria & INT & NOT NULL, FK & Maquinaria controlada (ON DELETE CASCADE). \\
fecha & DATE & NOT NULL & Fecha de la jornada de trabajo. \\
hora\_salida & TIME & NOT NULL & Hora de inicio o salida a campo. \\
hora\_devolucion & TIME & NULL & Hora de retorno o aparcamiento final. \\
\bottomrule
\end{tabularx}
\caption{Estructura técnica de la tabla \texttt{asistencia\_maquinaria}.}
\label{tab:asistencia_maquinaria}
\end{table}

% 13. registros
\subsection{Tabla \texttt{registros}}
Partes diarios de control horario y asistencia de personal obrero.
\begin{table}[H]
\centering\small
\renewcommand{\arraystretch}{1.3}
\begin{tabularx}{\textwidth}{>{\ttfamily\bfseries\color{midnightnavy}\raggedright\arraybackslash}p{2.7cm} >{\ttfamily\raggedright\arraybackslash}p{3.5cm} >{\raggedright\arraybackslash}p{2.2cm} Y}
\rowcolor{slate!8}
\textbf{Campo} & \textbf{Tipo de Dato} & \textbf{Nulo / Default} & \textbf{Descripción y Reglas de Integridad} \\
\specialrule{1pt}{0pt}{2pt}
id\_registro & INT AUTO\_INCREMENT & NOT NULL, PK & Identificador del parte de asistencia. \\
fecha & DATE & NOT NULL & Fecha de la jornada laboral registrada. \\
hora\_entrada & TIME & NOT NULL & Hora de fichada matutina o ingreso. \\
hora\_salida & TIME & NOT NULL & Hora de egreso o cierre de la jornada. \\
horas\_trabajadas & DECIMAL(5,2) & NULL & Cómputo neto de horas operativas. \\
id\_obrero & INT & NOT NULL, FK & Obrero que prestó servicios (CASCADE). \\
id\_obra & INT & NOT NULL, FK & Obra de asignación (CASCADE). \\
id\_usuario & INT & NULL, FK & Capataz certificante (ON DELETE SET NULL). \\
\bottomrule
\end{tabularx}
\caption{Estructura técnica de la tabla \texttt{registros}.}
\label{tab:registros}
\end{table}

% 14. recursos
\subsection{Tabla \texttt{recursos}}
Insumos, materiales y recursos extraordinarios consumidos en la jornada.
\begin{table}[H]
\centering\small
\renewcommand{\arraystretch}{1.3}
\begin{tabularx}{\textwidth}{>{\ttfamily\bfseries\color{midnightnavy}\raggedright\arraybackslash}p{2.7cm} >{\ttfamily\raggedright\arraybackslash}p{3.5cm} >{\raggedright\arraybackslash}p{2.2cm} Y}
\rowcolor{slate!8}
\textbf{Campo} & \textbf{Tipo de Dato} & \textbf{Nulo / Default} & \textbf{Descripción y Reglas de Integridad} \\
\specialrule{1pt}{0pt}{2pt}
id\_recurso & INT AUTO\_INCREMENT & NOT NULL, PK & Identificador del insumo consumido. \\
id\_obra & INT & NOT NULL, FK & Obra a la que se imputa el gasto (CASCADE). \\
id\_registro & INT & NULL, FK & Parte diario vinculado (ON DELETE SET NULL). \\
fecha & DATE & NOT NULL & Fecha del suministro en obra. \\
nombre & VARCHAR(150) & NOT NULL & Denominación del material o servicio. \\
cantidad & DECIMAL(10,2) & NOT NULL & Unidades o volumen empleado. \\
precio\_unitario & DECIMAL(10,2) & NULL & Costo unitario al momento del registro. \\
es\_material & BOOLEAN & DEFAULT 0 & Flag indicador (1: Material físico, 0: Servicio). \\
\bottomrule
\end{tabularx}
\caption{Estructura técnica de la tabla \texttt{recursos}.}
\label{tab:recursos}
\end{table}

% 15. combustible
\subsection{Tabla \texttt{combustible}}
Control de abastecimiento energético y consumo de carburantes en obra.
\begin{table}[H]
\centering\small
\renewcommand{\arraystretch}{1.3}
\begin{tabularx}{\textwidth}{>{\ttfamily\bfseries\color{midnightnavy}\raggedright\arraybackslash}p{2.7cm} >{\ttfamily\raggedright\arraybackslash}p{3.5cm} >{\raggedright\arraybackslash}p{2.2cm} Y}
\rowcolor{slate!8}
\textbf{Campo} & \textbf{Tipo de Dato} & \textbf{Nulo / Default} & \textbf{Descripción y Reglas de Integridad} \\
\specialrule{1pt}{0pt}{2pt}
id\_combustible & INT AUTO\_INCREMENT & NOT NULL, PK & Identificador de la carga de carburante. \\
nombre\_combustible & VARCHAR(100) & NOT NULL & Tipo de carburante (Diesel 500, Premium). \\
litros & DECIMAL(10,2) & DEFAULT 0.00 & Volumen suministrado en litros. \\
precio\_unitario & DECIMAL(10,2) & DEFAULT 0.00 & Costo por litro vigente a la fecha de carga. \\
precio\_total & DECIMAL(12,2) & DEFAULT 0.00 & Importe total de la operación monetaria. \\
fecha & DATE & NOT NULL & Fecha del suministro. \\
id\_obra & INT & NOT NULL, FK & Obra imputable del gasto (ON DELETE CASCADE). \\
id\_maquinaria & INT & NULL, FK & Máquina abastecida (ON DELETE SET NULL). \\
id\_factura & INT & NULL & Número de remito fiscal o comprobante. \\
\bottomrule
\end{tabularx}
\caption{Estructura técnica de la tabla \texttt{combustible}.}
\label{tab:combustible}
\end{table}

% 16. costos_generales
\subsection{Tabla \texttt{costos\_generales}}
Erogaciones periódicas fijas y variables de administración y soporte de obras.
\begin{table}[H]
\centering\small
\renewcommand{\arraystretch}{1.3}
\begin{tabularx}{\textwidth}{>{\ttfamily\bfseries\color{midnightnavy}\raggedright\arraybackslash}p{2.7cm} >{\ttfamily\raggedright\arraybackslash}p{3.5cm} >{\raggedright\arraybackslash}p{2.2cm} Y}
\rowcolor{slate!8}
\textbf{Campo} & \textbf{Tipo de Dato} & \textbf{Nulo / Default} & \textbf{Descripción y Reglas de Integridad} \\
\specialrule{1pt}{0pt}{2pt}
id\_costo\_general & INT AUTO\_INCREMENT & NOT NULL, PK & Identificador del costo operativo. \\
concepto & VARCHAR(150) & NOT NULL & Motivo de la erogación (alquiler, fletes, etc.). \\
categoria & ENUM('Fijo', 'Variable') & NOT NULL & Clasificación de contabilidad de gestión. \\
periodo & DATE & NOT NULL & Mes o período fiscal de imputación. \\
monto & DECIMAL(12,2) & NOT NULL & Importe monetario del gasto. \\
fecha\_registro & DATE & NOT NULL & Fecha de carga contable. \\
id\_usuario & INT & NOT NULL, FK & Usuario que autorizó la erogación. \\
\bottomrule
\end{tabularx}
\caption{Estructura técnica de la tabla \texttt{costos\_generales}.}
\label{tab:costos_generales}
\end{table}

% 17. partes_diarios
\subsection{Tabla \texttt{partes\_diarios}}
Reportes consolidados de operación y horas efectivas de maquinaria en obra.
\begin{table}[H]
\centering\small
\renewcommand{\arraystretch}{1.3}
\begin{tabularx}{\textwidth}{>{\ttfamily\bfseries\color{midnightnavy}\raggedright\arraybackslash}p{2.7cm} >{\ttfamily\raggedright\arraybackslash}p{3.5cm} >{\raggedright\arraybackslash}p{2.2cm} Y}
\rowcolor{slate!8}
\textbf{Campo} & \textbf{Tipo de Dato} & \textbf{Nulo / Default} & \textbf{Descripción y Reglas de Integridad} \\
\specialrule{1pt}{0pt}{2pt}
id\_parte & INT AUTO\_INCREMENT & NOT NULL, PK & Identificador del parte operativo. \\
id\_obra & INT & NOT NULL, FK & Obra donde operó el equipo (CASCADE). \\
id\_maquinaria & INT & NOT NULL, FK & Maquinaria controlada (CASCADE). \\
id\_usuario & INT & NOT NULL, FK & Operador o capataz certificante. \\
fecha & DATE & NOT NULL & Fecha del turno de trabajo. \\
horas\_trabajadas & DECIMAL(5,2) & DEFAULT 0.00 & Horas reloj con motor en operación activa. \\
horas\_paradas & DECIMAL(5,2) & DEFAULT 0.00 & Horas de inactividad por avería o lluvia. \\
litros\_combustible & DECIMAL(10,2) & DEFAULT 0.00 & Consumo estimado en la jornada. \\
observaciones & TEXT & NULL & Desperfectos mecánicos o incidencias. \\
\bottomrule
\end{tabularx}
\caption{Estructura técnica de la tabla \texttt{partes\_diarios}.}
\label{tab:partes_diarios}
\end{table}

% 18. notificaciones_leidas
\subsection{Tabla \texttt{notificaciones\_leidas}}
Marcado individual de alertas leídas por cada usuario del sistema.
\begin{table}[H]
\centering\small
\renewcommand{\arraystretch}{1.3}
\begin{tabularx}{\textwidth}{>{\ttfamily\bfseries\color{midnightnavy}\raggedright\arraybackslash}p{2.7cm} >{\ttfamily\raggedright\arraybackslash}p{3.5cm} >{\raggedright\arraybackslash}p{2.2cm} Y}
\rowcolor{slate!8}
\textbf{Campo} & \textbf{Tipo de Dato} & \textbf{Nulo / Default} & \textbf{Descripción y Reglas de Integridad} \\
\specialrule{1pt}{0pt}{2pt}
id\_usuario & INT & NOT NULL, PK, FK & Usuario que interactuó con la alerta (CASCADE). \\
tipo & VARCHAR(20) & NOT NULL, PK & 'maquinaria' (certificado) o 'obrero' (contrato). \\
id\_referencia & INT & NOT NULL, PK & ID del certificado o contrato descartado. \\
fecha\_leido & DATETIME & CURRENT\_TIMESTAMP & Marca temporal del acuse de recibo. \\
\bottomrule
\end{tabularx}
\caption{Estructura técnica de la tabla \texttt{notificaciones\_leidas}.}
\label{tab:notificaciones_leidas}
\end{table}

% ==============================================================================
% ==============================================================================
% CAPÍTULO 5: ESPECIFICACIÓN DETALLADA DE MÓDULOS
% ==============================================================================
\chapter{Especificación Detallada de Módulos Operativos}

\section{Módulo 1: Autenticación, Sesión y Control RBAC}

El subsistema de autenticación protege todos los puntos de entrada mediante validación criptográfica y control de tasa en ambos runtimes:
\begin{itemize}[leftmargin=1.5em]
    \item \textbf{Cifrado Cruzado BCrypt:} El sistema utiliza el algoritmo adaptativo BCrypt con costo computacional estándar. Tanto PHP (\texttt{password\_hash}) como C\# (\texttt{BCrypt.Net.BCrypt.Verify}) son 100\% interoperables, procesando indistintamente prefijos de formato \texttt{\$2y\$} y \texttt{\$2a\$}.
    \item \textbf{Control de Fuerza Bruta en Base de Datos:} En cada intento fallido, se registra la tupla en la tabla \texttt{intentos\_login}. Antes de autenticar, el sistema valida que la cuenta no acumule $\ge 5$ fallos o la IP $\ge 20$ en los últimos 5 minutos (código HTTP 429).
    \item \textbf{Filtrado RBAC en el Cliente (\texttt{auth.js}):} El script oculta de forma preventiva el DOM (\texttt{document.documentElement.style.visibility = 'hidden'}) hasta validar la sesión, evitando el parpadeo de contenido desprotegido (FOUC) a usuarios no autorizados.
\end{itemize}

\begin{apibox}{api/login.php}
    \textbf{Método HTTP:} \texttt{POST}\\[0.1cm]
    \textbf{Cuerpo JSON:} \texttt{\{"usuario": "admin", "password": "..."\}}\\[0.1cm]
    \textbf{Respuesta:} \texttt{\{"success": true, "user": \{"nombre": "...", "rol": "Administrador"\}\}}
\end{apibox}

\begin{methodbox}{HandleLoginAsync(JsonObject payload, string requestId)}
    Valida credenciales con \texttt{BCrypt.Net-Next}, inicializa la sesión en memoria en el host nativo de Windows y responde a la vista mediante \texttt{PostWebMessageAsStringAsync}.
\end{methodbox}

\electricsep

\section{Módulo 2: Obras Civiles, Contratos y Desglose Presupuestario}

Administra los proyectos constructivos y su trazabilidad financiera. Permite crear proyectos, asociar documentación contractual (almacenada en \texttt{LONGBLOB}), desglosar partidas en \texttt{contrato\_tareas}, y efectuar el cierre formal de contratos.

\begin{tipbox}{TRANSACCIONALIDAD EN EL CIERRE DE CONTRATOS}
    Al cerrar un contrato (por nuevo contrato o finalización de obra), tanto el endpoint PHP como el método C\# deben envolver la operación en una transacción ACID:
    \begin{enumerate}[leftmargin=1.2em]
        \item Calcular la sumatoria de tareas completadas: $\sum(\text{importe de tareas completadas})$.
        \item Actualizar \texttt{monto\_liquidado} e \texttt{importe\_final} en la tabla \texttt{contratos}.
        \item Cambiar el estado del contrato a \texttt{'Cerrado'}.
        \item Si se renueva con una adenda, crear el nuevo contrato referenciando \texttt{id\_contrato\_origen}.
    \end{enumerate}
\end{tipbox}

\begin{apibox}{api/obras.php}
    \textbf{Métodos Soportados:} GET (listar obras con saldo de contrato activo), POST (creación, edición y cierre de contratos con carga de archivos adjuntos).
\end{apibox}

\begin{methodbox}{HandleObrasListarAsync / HandleGuardarObraAsync}
    Ejecuta transacciones directas mediante ADO.NET con \texttt{MySqlTransaction}, garantizando consistencia relacional atómica.
\end{methodbox}

\electricsep

\section{Módulo 3: Cuadrilla de Obreros, Legajos y Contratos Laborales}

Gestiona el padrón del personal operativo y el control riguroso de vencimientos de contratos laborales.
\begin{itemize}[leftmargin=1.5em]
    \item \textbf{Padrón y Documentación:} Cada operario cuenta con registro de DNI/CUIT único, oficio especializado (\texttt{cargo}) y legajo digital en \texttt{contrato\_obrero}.
    \item \textbf{Alertas de Caducidad:} La vista de inicio detecta contratos cuya fecha de vencimiento ocurra dentro de los próximos 7 o 15 días, disparando notificaciones visuales para prevenir que un trabajador ingrese a obra sin seguro o contrato vigente.
\end{itemize}

\begin{apibox}{api/obreros.php}
    \textbf{Rutas y Operaciones Soportadas:}
    \begin{itemize}[leftmargin=1.2em, topsep=2pt, itemsep=1pt]
        \item \texttt{GET api/obreros.php}: Padrón general con estado de contrato.
        \item \texttt{POST api/obreros.php?action=guardar}: Alta y modificación de operarios.
        \item \texttt{POST api/obreros.php?action=subir\_contrato}: Carga de documentación digital.
    \end{itemize}
\end{apibox}

\begin{methodbox}{HandleObrerosListarAsync / HandleObrerosSubirContrato}
    Inserta o actualiza tuplas en \texttt{obreros} y \texttt{contrato\_obrero} con streaming de bytes de archivo.
\end{methodbox}

\electricsep

\section{Módulo 4: Parque de Maquinaria y Certificaciones Técnicas}

Controla los activos rodantes y maquinaria pesada afectada a obras constructivas:
\begin{itemize}[leftmargin=1.5em]
    \item \textbf{Gestión de Flota:} Registro de marca, modelo y estado operativo en \texttt{maquinaria}.
    \item \textbf{Certificados y Habilitaciones:} Pólizas de seguro, inspecciones técnicas vehiculares y habilitaciones de seguridad con fecha de caducidad en la tabla \texttt{certificado}.
    \item \textbf{Asignación a Obras:} Movimientos de entrada y salida asentados en \texttt{obra\_maquinaria} y \texttt{asistencia\_maquinaria}.
\end{itemize}

\begin{apibox}{api/maquinaria.php / api/cert\_maq.php}
    Administración de equipos mecánicos y subida de archivos binarios de certificación.
\end{apibox}

\begin{methodbox}{HandleMaquinariaListar / HandleMaquinariaCertSubir}
    Lectura de flota con verificación de vencimientos y persistencia de certificados en MySQL.
\end{methodbox}

\electricsep

\section{Módulo 5: Partes Diarios de Asistencia, Recursos y Combustible}

Constituye el núcleo transaccional cotidiano de las operaciones en obra:
\begin{itemize}[leftmargin=1.5em]
    \item \textbf{Control de Asistencia Horaria (\texttt{registros}):} Permite asentar hora de entrada y salida por obrero, computando automáticamente horas netas trabajadas.
    \item \textbf{Insumos y Materiales (\texttt{recursos}):} Asienta el uso de materiales de construcción o contratación de servicios especiales por jornada.
    \item \textbf{Suministro de Combustible (\texttt{combustible}):} Imputa cargas de gasoil o nafta directamente a la obra y máquina receptora, calculando el importe total monetario.
\end{itemize}

\begin{apibox}{api/guardar\_asistencia.php / api/guardar\_combustible.php}
    Recepción de partes diarios por lote y registro de consumo de combustible.
\end{apibox}

\begin{methodbox}{HandleGuardarAsistencia / HandleGuardarCombustible}
    Despacho transaccional de partes diarios en C\# con validación de no duplicidad de jornada por operario.
\end{methodbox}

\electricsep

\section{Módulo 6: Dashboard, Analíticas y Notificaciones Persistentes}

El panel de inicio consolida los indicadores clave de rendimiento (KPIs) de la constructora:
\begin{itemize}[leftmargin=1.5em]
    \item Conteo activo de proyectos en ejecución y presupuestos consolidados.
    \item Personal obrero presente en obra durante la jornada y cuadrillas asignadas.
    \item Sistema de notificaciones no destructivas: Las alertas de vencimiento descartadas se asientan en \texttt{notificaciones\_leidas} con clave compuesta \texttt{(id\_usuario, tipo, id\_referencia)}, evitando alertas repetitivas para ese usuario específico.
\end{itemize}

\begin{apibox}{api/dashboard.php / api/marcar\_notificacion.php}
    Consulta agregada de analíticas y persistencia de acuses de recibo de alertas.
\end{apibox}

\begin{methodbox}{HandleDashboardMetricas / HandleNotificacionMarcar}
    Cálculo de métricas consolidadas en memoria y persistencia de alertas leídas.
\end{methodbox}

\electricsep

\section{Módulo 7: Consultas Parametrizadas y Reportabilidad}

Permite la auditoría transversal y generación de informes de actividad operativa:
\begin{itemize}[leftmargin=1.5em]
    \item Filtro multifactorial por rango de fechas, obra específica y operario.
    \item Cómputo agregado de horas hombre trabajadas y costo acumulado de combustible.
    \item Vista tabular optimizada con paginación en cliente.
\end{itemize}

\begin{apibox}{api/consultas.php}
    \textbf{Parámetros GET:} \texttt{id\_obra}, \texttt{id\_obrero}, \texttt{fecha\_desde}, \texttt{fecha\_hasta}.
\end{apibox}

\begin{methodbox}{HandleConsultasBuscarAsync}
    Construcción dinámica y segura de consultas SQL preparadas con parámetros vinculados.
\end{methodbox}

\electricsep

\section{Módulo 8: Gestión de Usuarios y Auditoría Forense}

Administración del control de acceso y fiscalización de eventos del sistema:
\begin{itemize}[leftmargin=1.5em]
    \item \textbf{Gobernanza de Cuentas:} Creación y edición de usuarios con roles \texttt{Administrador} o \texttt{Capataz}.
    \item \textbf{Regla del Último Administrador:} Se impide la desactivación o eliminación de la última cuenta con privilegios de administrador en el sistema.
    \item \textbf{Bitácora Forense:} Registro automático de cada mutación en la tabla inmutable \texttt{auditoria\_logs}.
\end{itemize}

\begin{apibox}{api/usuarios.php / api/auditoria.php}
    CRUD de usuarios y consulta paginada de registros históricos de auditoría.
\end{apibox}

\begin{methodbox}{HandleUsersListAsync / HandleUserSaveAsync}
    Administración de operadores en C\# con hashing seguro y registro de eventos.
\end{methodbox}

\electricsep

\section{Módulo 9: Exportación Masiva de Datos a Formato Plano (CSV/ZIP)}

Permite la extracción completa de información operativa para conciliaciones contables externas:
\begin{itemize}[leftmargin=1.5em]
    \item Generación de flujos de texto plano en formato CSV con delimitador estandarizado (punto y coma o coma) y codificación UTF-8 con BOM para compatibilidad total con Microsoft Excel.
    \item Descarga empaquetada en archivos comprimidos ZIP conteniendo contratos y reportes consolidados por período.
\end{itemize}

\begin{apibox}{api/exportar\_csv.php}
    Descarga directa mediante streaming HTTP con cabecera \texttt{Content-Disposition: attachment; filename=export.csv}.
\end{apibox}

% ==============================================================================
\chapter{Protocolo de Comunicación Inter-Procesos (IPC en WebView2)}

\section{Mecanismo de Despacho y Manejo de Concurrencia}

En el modo de escritorio, la ventana de Windows Forms no realiza conexiones HTTP locales. Toda la comunicación fluye a través del canal IPC nativo de Microsoft Edge WebView2:

\begin{figure}[H]
    \centering
    \noindent\resizebox{\linewidth}{!}{%
    \begin{tikzpicture}[
        node distance=2.5cm and 3cm,
        proc/.style={rectangle, rounded corners=1.5mm, draw=midnightnavy, line width=1.2pt, fill=white, inner sep=7pt, align=center, font=\small},
        flow/.style={-{Stealth[scale=1.1]}, line width=1.3pt, color=electricblue}
    ]
        \node[proc, fill=subtleblue, text width=4.5cm] (js) {
            \textbf{\color{midnightnavy}JavaScript (Frontend)}\\[0.1cm]
            \footnotesize 1. Genera \texttt{requestId}\\
            2. Registra listener temporal\\
            3. Invoca \texttt{postMessage(json)}
        };

        \node[proc, right=3.5cm of js, fill=indigo!5!white, text width=5cm] (cs) {
            \textbf{\color{midnightnavy}C\# MainWindow.cs}\\[0.1cm]
            \footnotesize 1. Recibe en \texttt{OnWebMessageReceived}\\
            2. Despacha por \texttt{messageType}\\
            3. Ejecuta comando MySQL tipado\\
            4. Responde con \texttt{PostWebMessage}
        };

        \node[proc, below=1.8cm of js, fill=green!5!white, text width=4.5cm] (resolve) {
            \textbf{\color{midnightnavy}Resolución en JS}\\[0.1cm]
            \footnotesize 1. Filtra por \texttt{requestId}\\
            2. Remueve listener temporal\\
            3. Resuelve \texttt{Promise (resolve/reject)}
        };

        \draw[flow] (js.east) -- node[above, font=\scriptsize\color{charcoal}] {Mensaje JSON IPC} (cs.west);
        \draw[flow] (cs.south) |- node[below, font=\scriptsize\color{charcoal}] {Respuesta JSON con \texttt{requestId}} (resolve.east);
        \draw[flow] (resolve.north) -- node[left, font=\scriptsize\color{charcoal}] {Actualiza UI} (js.south);
    \end{tikzpicture}%
    }
    \caption{Ciclo de Vida de una Petición IPC Bidireccional en WebView2.}
    \label{fig:ipc_cycle}
\end{figure}

\begin{lstlisting}[language=javascript, caption={Función Universal de Despacho IPC en el Cliente (dashboard.js)}]
function sendDesktopRequest(type, payload, responseType) {
    return new Promise((resolve, reject) => {
        if (!window.chrome?.webview) {
            reject(new Error("No se detecto el entorno de escritorio."));
            return;
        }
        const requestId = `${type}_${Date.now()}_${Math.random().toString(16).slice(2)}`;
        const timer = setTimeout(() => {
            window.chrome.webview.removeEventListener("message", onMessage);
            reject(new Error("Tiempo de espera agotado en host C#."));
        }, 12000);

        function onMessage(event) {
            let data = event.data;
            if (typeof data === "string") {
                try { data = JSON.parse(data); } catch { return; }
            }
            if (!data || data.requestId !== requestId) return;
            window.chrome.webview.removeEventListener("message", onMessage);
            clearTimeout(timer);
            if (data.success) resolve(data);
            else reject(new Error(data.error || "Error al procesar la solicitud."));
        }

        window.chrome.webview.addEventListener("message", onMessage);
        window.chrome.webview.postMessage(JSON.stringify({ type, requestId, ...payload }));
    });
}
\end{lstlisting}

% ==============================================================================
% CAPÍTULO 7: MATRIZ DE REFERENCIAS CRUZADAS
% ==============================================================================
\chapter{Matriz de Referencias Cruzadas y Mapeo de Métodos}

\section{Mapeo Completo de Componentes de Frontend, Backend y Persistencia}

La Tabla \ref{tab:matriz_completa} constituye la guía definitiva para rastrear cualquier flujo funcional del sistema:

Para erradicar ambigüedades de mantenimiento y prevenir desbordamientos en la visualización, la matriz de métodos se presenta estructurada por dominios funcionales con anchos de celda balanceados:

% Subtabla 1: Seguridad y Auditoría
\subsection{Mapeo: Seguridad, Sesión y Auditoría}
\begin{table}[H]
\centering\small
\renewcommand{\arraystretch}{1.3}
\begin{tabularx}{\textwidth}{>{\bfseries\color{midnightnavy}\raggedright\arraybackslash}p{2.6cm} >{\raggedright\arraybackslash}p{2.5cm} >{\ttfamily\raggedright\arraybackslash}p{3.2cm} >{\ttfamily\raggedright\arraybackslash}p{3.4cm} Y}
\rowcolor{slate!8}
\textbf{Acción} & \textbf{Vista UI} & \textbf{Endpoint PHP} & \textbf{Despachador C\#} & \textbf{Tablas Afectadas} \\
\specialrule{1pt}{0pt}{2pt}
Autenticación & \texttt{Login.html} & \texttt{api/login.php} & \texttt{HandleLoginAsync} & \texttt{usuarios, intentos\_login} \\
Cierre Sesión & Navbar UI & \texttt{api/logout.php} & \texttt{HandleLogout} & \texttt{auditoria\_logs} \\
Gestión Usuarios & \texttt{Usuarios.html} & \texttt{api/usuarios.php} & \texttt{HandleUsersListAsync} & \texttt{usuarios, auditoria\_logs} \\
Auditoría Logs & \texttt{Auditoria.html} & \texttt{api/auditoria.php} & \textit{(Vía Web API)} & \texttt{auditoria\_logs} \\
\bottomrule
\end{tabularx}
\caption{Correspondencia funcional: Subsistema de Seguridad y Auditoría.}
\label{tab:matriz_seguridad}
\end{table}

% Subtabla 2: Obras y Contratos
\subsection{Mapeo: Obras Civiles, Contratos y Presupuesto}
\begin{table}[H]
\centering\small
\renewcommand{\arraystretch}{1.3}
\begin{tabularx}{\textwidth}{>{\bfseries\color{midnightnavy}\raggedright\arraybackslash}p{2.6cm} >{\raggedright\arraybackslash}p{2.5cm} >{\ttfamily\raggedright\arraybackslash}p{3.2cm} >{\ttfamily\raggedright\arraybackslash}p{3.4cm} Y}
\rowcolor{slate!8}
\textbf{Acción} & \textbf{Vista UI} & \textbf{Endpoint PHP} & \textbf{Despachador C\#} & \textbf{Tablas Afectadas} \\
\specialrule{1pt}{0pt}{2pt}
Listar Obras & \texttt{Obras.html} & \texttt{api/obras.php?GET} & \texttt{HandleObrasListarAsync} & \texttt{obras, contratos} \\
Detalle de Obra & \texttt{Obras.html} & \texttt{api/obras.php?id} & \texttt{HandleObrasDetalleAsync} & \texttt{obras, tareas, combustible} \\
Guardar Obra & \texttt{Obras.html} & \texttt{api/obras.php?POST}& \texttt{HandleGuardarObraAsync} & \texttt{obras, contratos} \\
Cerrar Contrato & \texttt{Obras.html} & \texttt{api/obras.php?cierre}& \texttt{HandleObrasCerrarContrato}& \texttt{contratos, contrato\_tareas} \\
\bottomrule
\end{tabularx}
\caption{Correspondencia funcional: Subsistema de Obras y Contratos.}
\label{tab:matriz_obras}
\end{table}

% Subtabla 3: Personal Obrero
\subsection{Mapeo: Personal Obrero y Legajos Laborales}
\begin{table}[H]
\centering\small
\renewcommand{\arraystretch}{1.3}
\begin{tabularx}{\textwidth}{>{\bfseries\color{midnightnavy}\raggedright\arraybackslash}p{2.6cm} >{\raggedright\arraybackslash}p{2.5cm} >{\ttfamily\raggedright\arraybackslash}p{3.2cm} >{\ttfamily\raggedright\arraybackslash}p{3.4cm} Y}
\rowcolor{slate!8}
\textbf{Acción} & \textbf{Vista UI} & \textbf{Endpoint PHP} & \textbf{Despachador C\#} & \textbf{Tablas Afectadas} \\
\specialrule{1pt}{0pt}{2pt}
Listar Obreros & \texttt{Obreros.html} & \texttt{api/obreros.php?GET} & \texttt{HandleObrerosListarAsync} & \texttt{obreros, contrato\_obrero} \\
Guardar Obrero & \texttt{Obreros.html} & \texttt{api/obreros.php?POST}& \texttt{HandleGuardarObreroAsync} & \texttt{obreros} \\
Subir Contrato & \texttt{Obreros.html} & \texttt{api/obreros.php?subir}& \texttt{HandleObrerosSubirContrato}& \texttt{contrato\_obrero} \\
Editar Fecha & \texttt{Obreros.html} & \texttt{api/obreros.php?edit}& \texttt{HandleObrerosEditarFecha}& \texttt{contrato\_obrero} \\
\bottomrule
\end{tabularx}
\caption{Correspondencia funcional: Subsistema de Obreros y Contratos.}
\label{tab:matriz_obreros}
\end{table}

% Subtabla 4: Maquinaria Pesada
\subsection{Mapeo: Maquinaria Pesada y Certificaciones}
\begin{table}[H]
\centering\small
\renewcommand{\arraystretch}{1.3}
\begin{tabularx}{\textwidth}{>{\bfseries\color{midnightnavy}\raggedright\arraybackslash}p{2.6cm} >{\raggedright\arraybackslash}p{2.5cm} >{\ttfamily\raggedright\arraybackslash}p{3.2cm} >{\ttfamily\raggedright\arraybackslash}p{3.4cm} Y}
\rowcolor{slate!8}
\textbf{Acción} & \textbf{Vista UI} & \textbf{Endpoint PHP} & \textbf{Despachador C\#} & \textbf{Tablas Afectadas} \\
\specialrule{1pt}{0pt}{2pt}
Listar Equipos & \texttt{Maquinaria.html}& \texttt{api/maquinaria.php} & \texttt{HandleMaquinariaListar} & \texttt{maquinaria, obra\_maq} \\
Guardar Equipo & \texttt{Maquinaria.html}& \texttt{api/maquinaria.php} & \texttt{HandleGuardarMaquinaria}& \texttt{maquinaria} \\
Listar Certs & \texttt{Maquinaria.html}& \texttt{api/cert\_maq.php} & \texttt{HandleMaquinariaCertListar}& \texttt{certificado} \\
Subir Cert & \texttt{Maquinaria.html}& \texttt{api/cert\_maq.php} & \texttt{HandleMaquinariaCertSubir}& \texttt{certificado} \\
\bottomrule
\end{tabularx}
\caption{Correspondencia funcional: Subsistema de Maquinaria y Certificados.}
\label{tab:matriz_maquinaria}
\end{table}

% Subtabla 5: Asistencia y Combustible
\subsection{Mapeo: Asistencias, Jornadas y Combustible}
\begin{table}[H]
\centering\small
\renewcommand{\arraystretch}{1.3}
\begin{tabularx}{\textwidth}{>{\bfseries\color{midnightnavy}\raggedright\arraybackslash}p{2.6cm} >{\raggedright\arraybackslash}p{2.5cm} >{\ttfamily\raggedright\arraybackslash}p{3.2cm} >{\ttfamily\raggedright\arraybackslash}p{3.4cm} Y}
\rowcolor{slate!8}
\textbf{Acción} & \textbf{Vista UI} & \textbf{Endpoint PHP} & \textbf{Despachador C\#} & \textbf{Tablas Afectadas} \\
\specialrule{1pt}{0pt}{2pt}
Guardar Asist & \texttt{Asistencia.html}& \texttt{api/guardar\_asist} & \texttt{HandleGuardarAsistencia}& \texttt{registros, recursos} \\
Cargar Catálogos& \texttt{Asistencia.html}& \texttt{api/datos\_asist} & \texttt{HandleAsistenciaCatalogos}& \texttt{obras, obreros, maq} \\
Carga Combustible&\texttt{Asistencia.html}& \texttt{api/guardar\_comb} & \texttt{HandleGuardarCombustible}& \texttt{combustible} \\
Consultas & \texttt{Consultas.html} & \texttt{api/consultas.php} & \texttt{HandleConsultasBuscarAsync}& \texttt{registros, obras, obreros} \\
\bottomrule
\end{tabularx}
\caption{Correspondencia funcional: Subsistema de Asistencia y Combustible.}
\label{tab:matriz_asistencia}
\end{table}

% Subtabla 6: Dashboard y Exportación
\subsection{Mapeo: Dashboard, Analíticas y Exportaciones}
\begin{table}[H]
\centering\small
\renewcommand{\arraystretch}{1.3}
\begin{tabularx}{\textwidth}{>{\bfseries\color{midnightnavy}\raggedright\arraybackslash}p{2.6cm} >{\raggedright\arraybackslash}p{2.5cm} >{\ttfamily\raggedright\arraybackslash}p{3.2cm} >{\ttfamily\raggedright\arraybackslash}p{3.4cm} Y}
\rowcolor{slate!8}
\textbf{Acción} & \textbf{Vista UI} & \textbf{Endpoint PHP} & \textbf{Despachador C\#} & \textbf{Tablas Afectadas} \\
\specialrule{1pt}{0pt}{2pt}
Métricas Resumen& \texttt{Inicio.html} & \texttt{api/dashboard.php} & \texttt{HandleDashboardMetricas}& Multi-tabla consolidada \\
Marcar Alerta & \texttt{Inicio.html} & \texttt{api/marcar\_notif} & \texttt{HandleNotificacionMarcar}& \texttt{notificaciones\_leidas} \\
Exportar CSV & \texttt{Exportar.html} & \texttt{api/exportar\_csv} & \textit{(Descarga directa stream)}& Tablas según entidad \\
\bottomrule
\end{tabularx}
\caption{Correspondencia funcional: Subsistema de Dashboard y Exportación.}
\label{tab:matriz_dashboard}
\end{table}

% ==============================================================================
% CAPÍTULO 8: CRIPTOGRAFÍA Y SEGURIDAD
% ==============================================================================
\chapter{Criptografía, Seguridad y Blindaje del Sistema}

\section{Defensa en Profundidad y Principios de Seguridad}

El sistema implementa defensas coordinadas en todas sus capas:
\begin{enumerate}
    \item \textbf{Cero Concatenación SQL (Inmunidad a SQLi):} En PHP, \texttt{PDO::ATTR\_EMULATE\_PREPARES = false} garantiza que el motor MySQL reciba el árbol de la consulta precompilado. En C\#, todos los parámetros se agregan mediante \texttt{MySqlParameter} tipado (\texttt{AddWithValue}).
    \item \textbf{Validación de Tipos MIME Reales:} Para evitar la subida de ejecutables camuflados, se analizan las firmas binarias (*Magic Bytes*) y se restringe a listas blancas estrictas (\texttt{pdf}, \texttt{doc}, \texttt{docx}, \texttt{jpg}, \texttt{jpeg}, \texttt{png}).
    \item \textbf{Protección contra Desactivación del Último Administrador:}
    Tanto \texttt{api/usuarios.php} como \texttt{MainWindow.HandleUserDeleteAsync} implementan la salvaguarda lógica:
    \begin{lstlisting}[language=SQL]
SELECT COUNT(*) FROM usuarios 
WHERE rol = 'Administrador' AND activo = 1 AND id_usuario <> @idUsuario;
    \end{lstlisting}
    Si el conteo arroja 0, la operación se cancela con error de negocio, evitando que el sistema quede huérfano de administradores.
\end{enumerate}

% ==============================================================================
% CAPÍTULO 9: PROCEDIMIENTOS DE MANTENIMIENTO
% ==============================================================================
\chapter{Guía Operativa de Mantenimiento y Procedimientos Paso a Paso}

\section{Procedimiento para Incorporar un Campo en una Entidad}

Para ejemplificar la disciplina de desarrollo requerida, se describe el proceso para añadir la columna \texttt{numero\_poliza} (\texttt{VARCHAR(100)}) a la entidad \texttt{maquinaria}:

\begin{enumerate}
    \item \textbf{Paso 1: Modificación de \texttt{schema.sql}:}
    \begin{lstlisting}[language=SQL]
ALTER TABLE maquinaria ADD COLUMN numero_poliza VARCHAR(100) NULL AFTER marca;
    \end{lstlisting}
    
    \item \textbf{Paso 2: Actualización de la API PHP (\texttt{api/maquinaria.php}):}
    \begin{lstlisting}[language=PHP]
// En responderListado():
$stmt = $pdo->query('SELECT id_maquinaria, nombre, marca, numero_poliza FROM maquinaria');

// En responderGuardado():
$poliza = limpiarTexto($body['numero_poliza'] ?? null, 100, false);
$stmt = $pdo->prepare('UPDATE maquinaria SET nombre = ?, marca = ?, numero_poliza = ? WHERE id_maquinaria = ?');
$stmt->execute([$nombre, $marca, $poliza, $id]);
    \end{lstlisting}

    \item \textbf{Paso 3: Actualización de C\# en \texttt{desktop-csharp/MainWindow.cs}:}
    \begin{lstlisting}[language=C]
// En HandleMaquinariaListarAsync:
numero_poliza = ReadNullableString(reader, "numero_poliza");

// En HandleGuardarMaquinariaAsync:
var poliza = SanitizeText(root, "numero_poliza", 100, required: false);
cmd.Parameters.AddWithValue("@numeroPoliza", ToDbValue(poliza));
    \end{lstlisting}

    \item \textbf{Paso 4: Actualización Visual en \texttt{web-ui/Maquinaria.html} y \texttt{maquinaria.js}:}
    Incorporar el campo \texttt{<input id="numeroPoliza">} en el modal de edición y enlazar el valor en el mapeo de renderizado y envío de formularios.
\end{enumerate}

\electricsep

\section{Compilación del Ejecutable de Escritorio (.NET 10)}

Para generar el ejecutable de producción para Windows:
\begin{lstlisting}[language=bash]
# Compilacion estandar en Release
dotnet build desktop-csharp/PorticoDesktop.csproj -c Release

# Publicacion autonoma de un solo archivo ejecutable (Self-Contained)
dotnet publish desktop-csharp/PorticoDesktop.csproj -c Release -r win-x64 --self-contained true -p:PublishSingleFile=true -p:IncludeNativeLibrariesForSelfExtract=true -o bin/Publish/
\end{lstlisting}
El binario final se ubicará en \texttt{bin/Publish/PorticoDesktop.exe}.

\electricsep

\section{Ejecución Local de Pruebas Automatizadas}

Pórtico cuenta con un test runner CLI sin dependencias externas en \texttt{tests/run.php}. Valida la integridad del esquema, la ausencia de registros huérfanos y la respuesta de los endpoints:
\begin{lstlisting}[language=bash]
php tests/run.php --verbose
\end{lstlisting}

% ==============================================================================
% CAPÍTULO 10: DIAGNÓSTICO Y RESOLUCIÓN DE PROBLEMAS
% ==============================================================================
\chapter{Diagnóstico y Resolución de Problemas Frecuentes (Troubleshooting)}

\section{Matriz de Diagnóstico y Árbol de Fallos}

\begin{table}[H]
\centering\small
\renewcommand{\arraystretch}{1.35}
\begin{tabularx}{\textwidth}{>{\bfseries\color{midnightnavy}\raggedright\arraybackslash}p{2.8cm} >{\raggedright\arraybackslash}p{3.2cm} >{\raggedright\arraybackslash}p{3.4cm} Y}
\rowcolor{slate!8}
\textbf{Síntoma} & \textbf{Causa Raíz Probable} & \textbf{Diagnóstico Rápido} & \textbf{Procedimiento de Reparación} \\
\specialrule{1pt}{0pt}{2pt}
Pantalla blanca en Host Desktop & Motor WebView2 no inicializado o fallo en esquema virtual. & Revisar si \texttt{WebView2 Runtime} está instalado en el SO. & Instalar Runtime Evergreen y validar \texttt{SetVirtualHostNameToFolderMapping}. \\
Error HTTP 403 CSRF Mismatch & Sesión expirada o pérdida de cabecera \texttt{X-CSRF-Token}. & Comprobar si la cookie de sesión fue purgada en el navegador. & Limpiar cookies o recargar la vista para renegociar el token criptográfico. \\
Carga truncada de archivo PDF & Archivo superior a \texttt{post\_max\_size} de PHP. & El payload llega vacío a \texttt{\$\_POST} y \texttt{\$\_FILES}. & Incrementar \texttt{upload\_max\_filesize = 50M} y \texttt{post\_max\_size = 55M} en \texttt{php.ini}. \\
Timeout en consultas C\# & Pool de conexiones MySQL agotado por bloqueos. & Revisar conexiones activas en MySQL (\texttt{SHOW PROCESSLIST}). & Asegurar el uso de \texttt{using (var conn = ...)} en todos los métodos C\#. \\
Falta de Paridad Web/Desktop & Método implementado en PHP pero omitido en C\#. & Cotejar la matriz de referencias cruzadas del Capítulo 7. & Implementar el mensaje en el bloque switch de \texttt{MainWindow.cs}. \\
\bottomrule
\end{tabularx}
\caption{Matriz de diagnóstico técnico y resolución de incidencias en producción.}
\label{tab:troubleshooting}
\end{table}

\electricsep

\section{Casos Frecuentes Detallados}

\begin{modernbox}{1. Error MySQL `errno 150` al importar esquema}
\textbf{Síntoma:} MySQL rechaza la creación de claves foráneas con el código de error 150.\\
\textbf{Causa:} Se intentó crear una referencia hacia una tabla aún no definida o con cotejamiento dispar.\\
\textbf{Solución:} El script DDL \texttt{schema.sql} debe abrir siempre con \texttt{SET FOREIGN\_KEY\_CHECKS = 0;} y cerrar con \texttt{SET FOREIGN\_KEY\_CHECKS = 1;}, garantizando además que todas las tablas declaren explícitamente \texttt{ENGINE=InnoDB DEFAULT CHARSET=utf8mb4 COLLATE=utf8mb4\_unicode\_ci}.
\end{modernbox}

\begin{modernbox}{2. Subida de contratos o planos falla silenciosamente en PHP}
\textbf{Síntoma:} La petición responde con error de campos vacíos o código HTTP 400.\\
\textbf{Causa:} El tamaño del archivo adjunto superó la directiva \texttt{post\_max\_size} del servidor PHP.\\
\textbf{Solución:} Configurar \texttt{.user.ini} o iniciar el servidor con directivas explícitas de memoria:
\texttt{php -d upload\_max\_filesize=25M -d post\_max\_size=25M -d memory\_limit=256M -S localhost:8000}.
\end{modernbox}

\begin{modernbox}{3. Pantalla en blanco en la Aplicación de Escritorio}
\textbf{Síntoma:} Al iniciar \texttt{PorticoDesktop.exe}, la ventana se abre pero no carga contenido.\\
\textbf{Causa:} Desfase en el nombre del recurso incrustado en el manifiesto binario de C\#.\\
\textbf{Solución:} Verificar que los archivos añadidos en \texttt{web-ui/} cumplan con la regla de inclusión en \texttt{PorticoDesktop.csproj}. La función \texttt{MapPathToResourceName} en \texttt{MainWindow.cs} transforma rutas separadas por barras en nombres lógicos separados por puntos (\texttt{webui.subcarpeta.archivo.ext}).
\end{modernbox}

\begin{modernbox}{4. Error HTTP 403: "Token de seguridad inválido"}
\textbf{Síntoma:} El formulario web rechaza peticiones de guardado tras periodos de inactividad.\\
\textbf{Causa:} La sesión de PHP expiró en el servidor pero la pestaña del navegador permaneció abierta.\\
\textbf{Solución:} Recargar la página para renegociar el token o solicitar una renovación asíncrona contra \texttt{api/csrf.php}.
\end{modernbox}

% ==============================================================================
% REFERENCIAS BIBLIOGRÁFICAS (AUTO-CONTENIDAS)
% ==============================================================================
% ==============================================================================
% REFERENCIAS BIBLIOGRÁFICAS (BIBTEX + ESTILO SBC)
% ==============================================================================
\nocite{*}
\bibliographystyle{sbc}
\bibliography{template}

\end{document}
